\chapter{The Module System}
\label{sec:module}


% Local Variables:
% mode: latex
% TeX-master: "CB-Manual"
% End:


Modules divide a {\em ConceptBase} database into
a hierarchy of namespaces that determine which objects are visible
inside the scope of a given module. Hence, each module forms a database that is part
of the whole database.
The scope of visibility also applies to  deductive rules, integrity constraints,
queries, active rules, and functions. They are also objects of the database and thus subject
to the visibility rules. Any object in the database belongs to exactly one module,
but it can also be visible to other modules, in particular to the sub-modules. 
Hence, whenever an object is created, updated or deleted (TELL, UNTELL), then
this has to be done in the context of the module that defines the object.
You can only query (ASK) an object, if the object is visible in your current
module context. Analogously, an object {\em x\/} can only reference other objects that are visible
in the module context where {\em x\/} is defined.

Modules that are not visible to each other also do not interfer with each other.
Hence, two users can use the same ConceptBase database and any change (TELL, UNTELL) done in 
the one module context does not influence the way how the other module reacts to requests.
Modules are organized in a hierarchy.
Changes to a supermodule shall impact all its sub-modules. This applies in particular to
the integrity of the sub-modules. For example, if an object in a sub-module refers to an object in
a super-module, then an update to the object in the super-module could render the sub-module
inconsistent. This would lead to a rejection of the update. Analogously, the result of
queries in a sub-module also depends on the objects in the super-modules. A module
in ConceptBase is quite similar to a {\tt DATABASE} in SQL in terms of isolating 
work spaces. On the other hand, the sub-module construct allows for controlled sharing of objects.

One useful application of modules in ConceptBase are modelling situations where different
objects are labeled with identical names. In earlier versions of {\em ConceptBase} this was
prohibited by the {\em Naming Axiom} which demands that different objects have different names.
Now objects with identical names can be stored in different modules without interferences.
Another application is to store alternative conceptual models of the same
domain in different modules. The alternative versions can also share a common core, e.g.
by storing it in a super-module of the version modules.

Another application of modules is to separate a metamodel (defining some constructs) from
the models represented in terms of the metamodel. The metamodel shall be defined in
a super-module and the models would be stored in sub-modules of this super-module. 
By this organization, the models do not interfer with each other unless they are
explicitely linked via export/import clauses. 
In multi-perspective modeling, models need to be rather tightly linked to each other.
Then, one should better store all such models in a single module. 

The set of visible objects in a module context is not limited  to the set of objects
defined for the module.
The {\em ConceptBase} module concept permits to reuse existing objects from different modules
via {\bf import} and {\bf export} interfaces. Furthermore,  modules can be {\bf nested}.
A nested module object (also called sub-module) is defined in the
context of another module, called its super modude.
A nested module object is only visible within the context of its super module.
Objects that are visible in the super module are visible (and can be reused) to all its
sub-modules. 

ConceptBase users can be assigned to individual {\bf home modules}, i.e. the module
they start working with when registering with the ConceptBase server. So-called
{\bf auto home modules} force every user in her own module to further reduce
potential unwanted interferences. A flexible {\bf access control} mechanism
allows to define access rights of users simply by query classes defined either globally
or locally to a module. 


\section{Definition of modules}

The class {\tt Module} defines an attribute {\tt contains} and thus is the construct to
create new modules. Each module is created by creating an instance of the class {\tt Module}.

\begin{verbatim}
Module in Class with
 attribute
    contains: Proposition
end
\end{verbatim}


Thus, a module is a container of objects. Modules create a name space: object names must be unique within
one module, but different modules can contain different objects with identical names.

{\bf Tell}, {\bf Untell} and {\bf Ask} operations work relatively to a specified module.
The normal way of specifying the context in which a transaction  takes place
is using the {\em Set Module} function of {\em CBIva}. See subsection \ref{sec:switching_modules}
to learn how to change the module context of a transaction. The command line interface
{\em CBShell\/} has similar operations to switch the module context.

The basic set of predefined objects of {\em ConceptBase} (such as {\tt Class}, {\tt Proposition}, {\tt QueryClass}, etc.) belongs to the predefined module {\tt System}. 
The default module of clients logging into a CBserver is {\tt oHome}, a direct submodule of {\tt System}.
You can set the module context to your other modules in order to manipulate them.

The {\tt contains} attribute of the {\tt Module} object is a derived attribute and is a link
to all objects defined inside the context of a certain module. This attribute is not stored explicitly
but can be used anywhere in the O-Telos assertion language.

Now we show how to create modules.
We introduce a small running example in order to demonstrate all module-related facilities of
{\em ConceptBase}.

Let's assume you've started ConceptBase with a fresh database.
After telling the following two frames, the server contains two new modules, which are nested inside the pre-defined {\tt oHome} module:


\begin{verbatim}
Master in Module end
Work in Module end
\end{verbatim}

The super-module {\tt oHome} includes the two objects {\tt Master} and {\tt Work} but not the
objects contained in them (compare also Figure \ref{fig:mod_hierarchy}).
We call {\tt oHome} the super-module of the sub-modules
{\tt Master} and {\tt Work}. One you switch to a sub-module, say {\tt Master}, then all
sub-sequent operations apply to this module context. Telling new objects makes
them members of this module context, in other words {\tt Master} contains them.


\section{Switching between module contexts} \label{sec:switching_modules}


The standard way for changing the module context is to choose the {\em select module} menu entry
from the {\em Options} menu in the user interface. A dialog with a listbox containing
all known modules in the current context are shown.
Double-clicking a module entry in the listbox sets the current module context to the
selected module and lists all modules that are visible in this module (i.e., its
submodules and supermodules.
\cbfigure{choose_modcontext}{5cm}{The module context selection window}

In our example, you should get a window displaying
the modules {\tt Master}, {\tt System}, {\tt Work}, and {\tt oHome} in alphabetic order
(see Figure \ref{fig:choose_modcontext}).
Now double-click the entry {\tt Work}. As a result the module context of
the {\em CBserver} is set to the module {\tt Work}. As the {\tt Work} module
has got no nested modules, there are no additional modules
displayed in the listbox (i.e., all modules that are visible in the {\tt oHome}
module are also visible in the {\tt Work} module.
Alternatively, you can specify the module context using the {\em load model\/} operation
and placing a

{\tt \{\$set module=MODNAME\}}

inline command placed before the first Telos frame
 of a  source model file.
{\tt MODNAME} stands for the name of the module in which you wish to define the
content of the source model file. Module paths like {\tt System/oHome/Work} are allowed as well.
Please note that only one inline-command is allowed within one
source model file.
Specifying module contexts using inline commands is a facility for automatically loading
Telos frames of large applications
which are spread over different modules -- without requiring the user to employ
the {\em select module\/} function from {\em CBIva}.

In CBShell, you can use the command {\tt cd} (alias: {\tt setModule}) to switch the module context.
Let's set the module context first to {\tt oHome} and then TELL the
following frames for a very simple ER notation (using CBShell syntax, see section \ref{cha:cbshell}):

\begin{verbatim}
cd oHome

tell "EntityType end
RelationshipType with
  attribute
     role: EntityType
end"
\end{verbatim}

Then set the module context first to {\tt Master} and then TELL the
Employee frame:

\begin{verbatim}
cd Master

tell "Employee in EntityType,Class with
   attribute
      name : String
   constraint
      nec: $ forall e/Employee exists n/String (e name n) $
end"
\end{verbatim}

The object {\tt Employee} is now only visible in the module {\tt Master}. When you
set the module context to {\tt Work} (or {\tt System}) and try to load
the {\tt Employee} object, you should get an error message from the server stating that
the object {\tt Employee} is not visible in that module context. The object 
{\tt Employee} correctly references {\tt EntityType}, since it is visible via
the super-module {\tt oHome}.

\section{Using nested modules}
\label{sec:nested_modules}

A nested module object is a module defined in the context of a super module and therefore
is contained in the super module. The objects contained in the nested module
are {\em not\/} contained in the super module.  

After the definition of the {\tt Master} and {\tt Work} modules as nested modules
to the {\tt oHome} module, let's define a nested module to the {\tt Work} module.
We assume that the current module context is set to the  {\tt Work} module.
Now tell the following module object:

\begin{verbatim}
cd Work
tell "Test in Module end"
\end{verbatim}

As a result we have defined the nesting hierarchy depicted in Figure \ref{fig:mod_hierarchy}.

\cbfigure{mod_hierarchy}{6.5cm}{Nested module hierarchy}

A nested module can see the content of all modules on the path to the root module {\tt System}.
Therefore, when you set the
module context to {\tt Test}, you can reference all objects contained
in the modules {\tt Test}, {\tt Work}, {\tt oHome}, and {\tt System}. When you set the module context
to {\tt Work}, you can reference all objects contained in {\tt Work}, {\tt oHome},
and {\tt System}.

Any ConceptBase database has the pre-defined modules {\tt System} and {\tt oHome}. The {\tt System}
module contains the pre-defined objects of ConceptBase, i.e. the O-Telos classes {\tt Proposition},
{\tt Individual}, {\tt Attribute}, {\tt InstanceOf}, and {\tt IsA}, but also
a large number of other pre-defined objects that are required for a functioning system,
e.g. {\tt Integer}, {\tt Class}, {\tt QueryClass}, and many more.
The {\tt oHome} module is initially empty. It is the default home module
for users and clients that add/update/delete and query objects. A user could delete objects in
the {\tt System} module, which could then disable core functions of ConceptBase. 
One can however also prevent such updates by limiting the access to the {\tt System}
module. 


\section{Exporting and importing objects}

In order to use objects which are not visible in a module, {\em ConceptBase} offers the
export/import facility.
The class {\tt Module} defines two further attributes, namely to specify objects exported from
a module and to modules imported by a module.

\begin{verbatim}
Module with
 attribute
    contains     : Proposition;
    exports      : Proposition;
    imports      : Module
end
\end{verbatim}

In order to allow other modules to import objects from a module, we need to define an {\tt export} attribute from the module object
to those objects. We call the set of
objects exported by a module the {\em export interface} of a module.
In order to include the export interface of another module to a module, we need to define
an {\tt imports}-attribute from the module object to the module to be imported.


In our running example we would like to define a specialization of the Class {\tt Employee} within
the  {\tt Work} module. It is desirable to reuse   {\tt Employee} from the
{\tt Master} module instead of redefining it in the  {\tt Work} module.

In order to import the class {\tt Employee} to the {\tt Work} module, we have to
define all objects belonging to the {\tt Employee} class as exported objects.
First change the module context to {\tt Master}. Now TELL the following frame
within module {\tt Master}:

\begin{verbatim}
Master with
   exports
      e1 : Employee;
      e2 : Employee!name
end
\end{verbatim}

Now you have defined the  objects {\tt Employee} and {\tt Employee!name}
as exported objects of the {\tt Master} module.
Any module in your database, which defines an {\tt imports}-attribute to the
{\tt Master} module, can now reference these  objects.
Now change the module context to {\tt Work} and TELL the following frame within
module {\tt Work} (see also Figure \ref{fig:mod_hierarchy}):

\begin{verbatim}
Work with
   imports
      i1 : Master
end
\end{verbatim}

The  objects mentioned above are visible in the context of the {\tt Work}
module. Check this by loading the {\tt Employee} object with the {\em Edit/Load Object} function
of {\em CBIva}.
The import declaration for {\tt Work} is done within the module context
{\tt Work}. The attribute {\tt Work!imports} is thus part of the {\tt Work}
module while the {\tt Work}  object is contained in the {\tt oHome} module.
The object {\tt Master} is visible in the {\tt Work} module as well.
The visibility rules restrict the possible import declarations. 
You can now define the specialization {\tt Manager} of {\tt Employee}  in the
context of the {\tt Work} module:


\begin{verbatim}
Manager isA Employee end
\end{verbatim}

When untelling {\tt exports} declarations from  a module,
{\em ConceptBase} checks for integrity violation in all concerned modules.
Try untelling the {\tt exports} attributes from the {\tt Master} module
and you should get an error message saying that referential integrity is violated
in the {\tt Work} module. The reason for this violation is simple: since the class
{\tt Employee} is no longer exported from {\tt Master}, it is no longer visible
inside the {\tt Work} module and therefore the referential integrity (a builtin O-Telos axiom)
is violated for the {\tt Manager} specialization of {\tt Employee}.

If you define integrity constraints in the exporting module {\tt Master}, then
these constraints are not checked for objects in the importing module {\tt Work}.
For example, the constraint {\tt nec} of {\tt Employee} is only visible in {\tt Master}
and its sub-modules, not in {\tt Work} and its sub-modules.
Hence, you can still declare an object like {\tt bill} in module {\tt Test} that does not need to
fulfill the constraint for {\tt Employee}: 

\begin{verbatim}
cd Test
tell "bill in Employee end"
\end{verbatim}

\section{Modules and metamodeling} \label{sec:metamodule}

The module hierarchy in Figure \ref{fig:mod_hierarchy} assigns objects belonging to different
abstraction levels (meta classes, classes, objects, ...) to different modules. This assignment
is not prescribed, but it is recommended. For example, the module {\tt oHome} stores
meta classes such as {\tt EntityType}. This metaclass is then used in the sub-module
{\tt Work} to define a class like {\tt Manager}. Finally, the module {\tt Test} instantiates
the class {\tt Manager} by {\tt bill}. There is a natural reason for such a structure. The
persons defining meta classes are engineering modeling languages. The persons defining 
classes are conceptual modelers or application programmers, and the persons defining objects
at the lowest abstraction level are application users. These activities depend on each other but
one should separate the workspaces to shield an update to the modeling language
from an update to a conceptual model. Another reason is that a modeling language 
can be used for many conceptual models. Each conceptual model needs the definitions of the
modeling language but they typically should be separated from each other. Hence, each conceptual
model can be stored in its own module, being a sub-module of the module defining the 
modeling language.

The following CBShell script shows how a module hierarchy is created.
The command "cd" is use to switch to a module, the command "mkdir" creates a new submodule
in the current module. The module {\tt oHome}
contains the submodule {\tt ERnotation} for the definition of the ER modeling language.
The two submodule {\tt UModel} and {\tt LibModel}. The submodule {\tt UData} of
{\tt UModel} stores a sample database for the university model.

\begin{verbatim}
cd oHome
mkdir ERnotation
cd ERnotation
tellModel ERD-Language.sml.txt
mkdir UModel
mkdir LibModel
cd ERnotation/UModel
tellModel UniversityModel.sml.txt
mkdir UData
cd ERnotation/UModel/UData
tellModel UniversityData.sml.txt
\end{verbatim}




\section{Setting user home modules} \label{sec:homemodule}
% written by M.Jeusfeld, 30-Jan-2004

When ConceptBase is used in a multi-user setting, it makes sense to
automatically assign clients of ConceptBase users to a dedicated
module context, their so-called {\em home module}.
To use this feature, the database of the ConceptBase server has to contain
instances of the pre-defined class {\tt CB\_User}. This class is defined as follows:

\begin{verbatim}
CB_User with
 attribute
      homeModule : Module
end
\end{verbatim}

Assume that we have two users {\tt mary}  and {\tt john} who need to
be assigned to different modules when they log into the CBserver by
their favorite user interface.
The system adminstrator should then include the following definitions
to the database of the CBserver:

\begin{verbatim}
Project1 in Module end
Project2 in Module end

mary in CB_User with
 homeModule m1 : Project1
end

john in CB_User with
 homeModule m1 : Project2
end
\end{verbatim}

As a consequence, the start module of the two users will be set accordingly when
they log into the CBserver. The home module feature is especially useful
in a teaching environment. The teacher can put some Telos models into
the shared {\tt oHome} module. Students' home modules would be assigned to sub-modules,
e.g. based on group membership. Each student group can then work on an assignment
by working on their sub-module without interfering with other student groups.

There is a subclass {\tt AutoHomeModule} of {\tt Module}, which supports
applications of ConceptBase where by default any user should
work in her own module context. Rather than having to define separate
modules for each user explicitely, you can just define a certain module to be an instance
of {\tt AutoHomeModule}. 

\begin{verbatim}
LectureModule in AutoHomeModule with
  exception e1: mary
end

mary in CB_User with
 homeModule m1 : LectureModule
end

john in CB_User with
  homeModule m1: LectureModule
end
\end{verbatim}

You can also define a rule to assign all or a subset of users to this 
module:

\begin{verbatim}
CB_User in Class with  
  rule
    homeRule : $ forall u/CB_User (u homeModule LectureModule) $
end 
\end{verbatim}


Here, user {\tt john} (a student) will be automatically be assigned to a new module
{\tt M\_john} that is created as sub module of {\tt LectureModule}.
User {\tt mary}, presumably a teacher, is defined to be an {\em exception\/} to this
rule and she will get the home module {\tt LectureModule}.
By this, one can reduce the chances of unwanted interferences between users
of the module {\tt LectureModule}. Still, all users can read the
definitions in the module {\tt LectureModule} and its submodules unless
access restrictions are defined (see section \ref{sec:modaccess}).

The simplest way to separate the workspaces of {\em any\/} user
is to tell

\begin{verbatim}
oHome in AutoHomeModule end
\end{verbatim}

In this case, no user needs to be defined explicitely%
\footnote{A side effect of the server method {\tt ENROLL\_ME} is that the user
of the registering client will automatically be defined as an instance of class
{\tt CB\_User}. The definition is be made in the context of module
{\tt oHome}.}
 as instance of {\tt CB\_User} and still
will get assigned her own sub module to work in. 
The auto-home module becomes active as soon as the module is declared
as instance of {\tt AutoHomeModule}. If you want to enable {\tt oHome}
as auto-home module from the very beginning when a database is created,
then you can activate it by a parameter of the CBserver, e.g.

\begin{verbatim}
cbserver -g public -new MYNEWDB
\end{verbatim}

This will instruct the CBserver to tell {\tt "oHome in AutoHomeModule end"}
when it sets up the new database. See also section \ref{sec:pubcbserver}.


The home module definitions need to be made within module {\tt oHome} because
they will be evaluated upon client registration (server method {\tt ENROLL\_ME})
in this module context.
Please note that the module context is only dependent on the user name,
not on the client and not on the network location of the user.
It could well be that a user {\tt mary} is defined on multiple
computers on the network and that different natural persons are
identified by {\tt mary}. ConceptBase currently cannot detect such cases.


\section{Limiting access to modules}
\label{sec:modaccess}
% written by M.Jeusfeld, 19-Feb-2004

When multiple users work on the same server, their workspaces not only need to be separated
in a controlled way by means of the module feature. Users are also interested in controlling who
has which rights on their workspace (=module). ConceptBase includes basic support for 
rights definition and enforcement via a user-definable query class {\tt CB\_Permitted}.
The signature of this query class has to conform the following format:

\begin{verbatim}
CB_Permitted in GenericQueryClass isA CB_User with
  parameter
    user: CB_User;
    res: Resource;
    op: CB_Operation
   ...
end
\end{verbatim}

A {\tt user} is allowed to perform the operation {\tt op} on the resource {\tt res}
iff the constraint of the query is satisfied. Then, {\tt user} is returned as answer
of the query. If not, the answer is {\tt nil} (equals empty set). 
Some fundamental definitions are pre-defined objects of ConceptBase:


\begin{verbatim}
Resource with end  
Module isA Resource end  
CB_Operation end   
CB_ReadOperation isA CB_Operation end  
CB_WriteOperation isA CB_Operation end
TELL in CB_WriteOperation end
ASK in CB_ReadOperation end   
\end{verbatim}

Hence, at least two operations {\tt TELL} and {\tt ASK} are pre-defined symbolizing write and read
accesses to a resource. Modules are the prime examples of resources to be access-protected.
Currently, only access to them is monitored by the CBserver.

When a user wants to switch to a new module, then he must at least have the permission
to execute the operation {\tt ASK} on it, i.e. read permission. Otherwise, the module
switch is rejected.
This check is the main protection scheme offered to module owners.
Define the query class {\tt CB\_Permitted} in the module that needs
protection. 

Assume that there is a user {\tt jonny} who wants to protect his module {\tt Mjonny}.
To do so, he would first define the module
 and set the module context to {\tt Mjonny}.

\begin{verbatim}
Mjonny in Module end
\end{verbatim}

Then, he would set the module context to {\tt Mjonny} define his rights management policy,
for example:

\begin{verbatim}
CB_Group with 
  attribute
    groupMember: CB_User;
    permitted_read: Resource;
    permitted_write: Resource;
    owner_resource: Resource
end

CB_User isA CB_Group end  

CB_Group in Class with
  rule
   rr1: $ forall p/Resource u/CB_User 
          (u owner_resource p) ==> (u permitted_write p) $;
   rr2: $ forall p/Resource u/CB_User 
          (u permitted_write p) ==> (u permitted_read p) $;
   rr3: $ forall u/CB_User (u groupMember u) $;  
   rr4: $ forall u/CB_User p/Resource
               ( exists g/CB_Group (g groupMember u) and 
                 (g owner_resource p) ) ==> (u owner_resource p) $;
   rr5: $ forall u/CB_User p/Resource
                ( exists g/CB_Group (g groupMember u) and 
                  (g permitted_write p) ) ==> (u permitted_write p) $;
   rr6: $ forall u/CB_User p/Resource
                ( exists g/CB_Group (g groupMember u) and 
                (g permitted_read p) ) ==> (u permitted_read p) $
end


CB_Permitted in GenericQueryClass isA CB_User with
  parameter
    user: CB_User;
    res: Resource;
    op: CB_Operation
  constraint
    cperm: $ (
              ( not exists u/CB_User 
                 (u owner_resource ~res) and 
                  not (u == ~user) ) 
              or
              ( (~op in CB_ReadOperation) and 
                (~user permitted_read ~res) ) 
                or
              ( (~op in CB_WriteOperation) and 
                (~user permitted_write ~res) )
             )
              and UNIFIES(~user,~this) $
end
\end{verbatim}

In the above example, access rights are granted to groups of ConceptBase users. 
The {\em owner} of a resource will always have full access via rules {\tt rr1} and {\tt rr2}.

Then, the user would claim ownership to the module via

\begin{verbatim}
jonny in CB_User with
  owner_resource r1: Mjonny
end
\end{verbatim}

Then, only {\tt jonny} can switch to the module {\tt Mjonny}. If a second user like {\tt mary}
was to be granted read permission, {\tt jonny} would define within module {\tt Mjonny}:

\begin{verbatim}
mary in CB_User with
  permitted_read r1: Mjonny
end
\end{verbatim}

In the above example, rights can also be granted to groups of users and then inherited to
its members via rules {\tt rr4} to {\tt rr6}.
It should be noted that the definitions of {\tt owner\_resource}, {\tt permitted\_read} and
{\tt permitted\_write} are just for illustrating what is possible. ConceptBase only
requires the query class {\tt CB\_Permitted} in the module where the access rights need
to be enforced. If such a query class (or a local version as explained below) is not
defined, then any access is permitted for any user.

When a user attempts to switch to new module context, ConceptBase checks whether
the user has at least read permission, i.e. permission for executing the
operation {\tt ASK}, on the module. If permission is not granted, the user 
cannot switch the module context and an error message is presented.


The definition of the query {\tt CB\_Permitted} is visible in the module where it is
defined and in all sub-modules of this module. One can also define a local version
of the query by appending the module name to its name, e.g. {\tt CB\_PermittedMjonny}.
This version is only tested for access to the module {\tt Mjonny}. The local
overrides the general version {\tt CB\_Permitted} and its function is not inherited to sub-modules. 
The subsequent definition prevents updates to the {\tt Test} module, if it is
visible in the {\tt Test} module and access control is enabled by the CBserver:

\begin{verbatim}
GenericQueryClass CB_PermittedTest isA CB_User with
  parameter
    user: CB_User;
    res: Resource;
    op: CB_Operation
  constraint
    cperm: $ (~op in CB_ReadOperation) and UNIFIES(~user,~this) $
end
\end{verbatim}


There are plenty of ways to combine general and local versions of {\tt CB\_Permitted}
yielding different access policies. When using access control, at least the module
{\tt System} should be protected. Otherwise, users could change essential definitions
affecting all other users. Examples for access control policies are in the HOW-TO section of the
ConceptBase Forum (\url{http://merkur.informatik.rwth-aachen.de/pub/bscw.cgi/2281940}).

It is very well possible to make access to a ConceptBase database completely impossible
by errors in the definition of {\tt CB\_Permitted}. For example, one could deny access
to any operation by the following simple rule:

\begin{verbatim}
CB_Permitted in GenericQueryClass isA CB_User with
  parameter
    user: CB_User;
    res: Resource;
    op: CB_Operation
  constraint
    cperm: $ FALSE $
end
\end{verbatim}

In such cases, one has to start the CBserver with disabled access control and repair the
definition of the query {\tt CB\_Permitted}. 
Access control is by default set to level 1, which just constrains the scope of an
UNTELL to the local module. You need to set the CBserver option {\tt -s} to 2 to
fully enable access control (see option -s in section \ref{cha:cbserver}).


\subsection{Access to System module} \label{sec:system_access}

The access definition via the query {\tt CB\_Permitted} or its localized variant
is available in the module {\tt oHome} and its submodules. It is not available
for checking permission to access the {\tt System} module. The reason is that
user details are stored in {\tt oHome} and signature of {\tt CB\_Permitted}
requires that the user details are known. The protection of the {\tt System}
module is instead configured via the security level of the CBserver (option -s):


\begin{itemize}
\item Level=0: No access control. Read and write operations are allowed to the
{\tt System} module.
\item Level=1: Read and writes are allowed but only if the current module is set
to {\tt System}.
\item Level=2: Read is permitted but write operations to {\tt System} are disallowed.
\end{itemize}



{\bf Caution:} The access control feature of ConceptBase avoids
some {\em unwanted\/} interferences in a setting where multiple users
work on the same server. The system in {\em not\/} save against malicious attacks!
Neither does it prevent all unwanted interferences.



\section{Listing the module content} \label{sec:module_content}
% written by M.Jeusfeld, 22-Feb-2011

The {\tt contains} attribute allows to check which objects belong to a module. It can be used
by a simple query that lists the module content of all modules that are currently visible:

\begin{verbatim}
ShowModule in QueryClass isA Module with
  computed_attribute
    cont : Proposition
  constraint
    ccont : $ (~this contains ~cont) $
end
\end{verbatim}

A more sophisticated method is to use the builtin query {\tt listModule}. A call of
{\tt listModule} without parameters will list the current module as Telos frames.
You can also provide the module to be listed as a parameter, e.g.

\begin{verbatim}
   listModule[System]
\end{verbatim}

will list the content of the module {\tt System}. ConceptBase will check
read permission before a module content is listed. You can also use
a module path as parameter, e.g.\ 

\begin{verbatim}
   listModule[System-oHome-Work]
\end{verbatim}

A module path is formed much like a directory path in a file system.
The root module is {\tt System} and modules names are separated by 
the character {\tt '-'}. You can also use '/' as module separator:
\begin{verbatim}
   listModule[System/oHome/Work]
\end{verbatim}

The implementation of the {\tt listModule} query preserves
the order in which objects have been created. Note that the {\tt System}
module is defining the essential objects that ConceptBase requires to
run correctly. You can list the {\tt System} module but you should
not change it. 

If the content of a module was created by separate transactions, then
{\tt listModule} shall indicate them by a separator line "{\tt \{---\}}".
This separator is disabled when you specify {\tt -mg whole}
as CBserver parameter (see section \ref{sec:cbsparams}).
The separator is technically a comment. However, CBGraph and CBIva
shall use such separator line to split a sequence
of frames into separate TELL transactions.
This is the default behaviour (or when you start the CBserver with
parameter {\tt -mg split}). The third option is to start the CBserver with
parameter {\tt -mg minsplit}. This minimizes the number of separator line.
They are only included if the previous transaction did create or modify a
deductive rule. The third option meant to spead up module loading
from sources. 

If a module path does not exist or the current user has no
read permission on modules in the path, then the answer
"\{* no *\}" is returned.

\subsection{Restrictions of {\tt listModule}} \label{sec:listmod}

The query {\tt listModule} extracts all objects of a module into a single
Telos source. Since a module is typically created by a sequence of
TELL/UNTELL transactions, this Telos source can in rare cases fail to
be told by a single transaction. An example is the specification of
the ERD model in 
\url{http://merkur.informatik.rwth-aachen.de/pub/bscw.cgi/188651}.
The university model (and then university data) depend on rules
defined in the model ERD-Semantics, specifically the semantics of
the {\tt ISA} type (instances of subclasses are also instances
of superclasses). If you tell the university model and the ERD semantics in
the same transaction, then the rules for the {\tt ISA} type are not yet 
usable for checking the consistency of the university model.
A way out is to store the university model as sub-module of 
the module that contains the ERD semantics, and university data
as sub-module of the university module.

A way out of this dilemma is to use the module feature of ConceptBase
in a thoughtful way. In this case, define a module like {\tt ERnotation}
to which the ERD notation is told. Within this module, define a submodule
like {\tt UModel}, to which the example ER model is told. Finally,
define a submodule {\tt UData} to hold the sample data. By nesting the
submodules in this way, the sample data objects can "see" the definitions
of the example ER model. And the example ER model can "see" the objects
of the ER notation.

\section{Purging a module} \label{sec:module_purge}
% written by M.Jeusfeld, 30-Oct-2018

The builtin query {\tt purgeModule} attempts to delete all propositions
of the current module or the module that is specified as a parameter.
The operation fails if the module contains a non-empty submodule. 
Further, the operation may non be applied to the {\tt System} module.

The operation requires that the user has appropriate permissions on the module
to be purged. 

The query {\tt purgeModule} is declared as a hidden object. Hence, it shall not show
up in the 'Display queries' dialogue of CBIva. The current implementation is experimental.
An alternative is to list the contents of the current module and then
applying the UNTELL operation to the whole content.





\section{Saving and loading module sources} \label{sec:module_sources}
% written by M.Jeusfeld, 22-Feb-2011

The ConceptBase server (see section \ref{cha:cbserver}) has two command line
options {\tt -save} and {\tt -load} to synchronize with the content of database modules as
Telos sources files in the file system of the CBserver. The purpose
of this feature is to allow an easy way to save/reload the complete content
of a database using a readable source format. These sources can be modified
with a regular text editor. 

The source files generated by the save function include a {\tt set module}
directive that instructs the CBserver to load the file to the original
module path when the load function is invoked (or when the file is loaded
manually via the {\tt load model} function of CBIva (see section \ref{cha:workbench}).
There are two ways to represent saved/loaded module sources in the file system:

\begin{description}
\item[flat directory structure:] If the module separator is set to '-' (see option -ms in
section \ref{sec:cbsparams}), the files names consist of a module path
starting with the root module name {\tt System} followed by the file type {\tt sml}.
For example, the file {\tt System-oHome-AB.sml} will hold the contents of
the module {\tt AB} that is a submodule to {\tt oHome} that is a submodule of {\tt System}.
\item[deep directory structure:] If the module separator is set to '/', then the module sources
are placed in sub-directories named like the modules.
For example, the file {\tt System/oHome/AB/AB.sml} will hold the contents of
the module {\tt AB} that is a submodule to {\tt oHome} that is a submodule of {\tt System}.

\end{description}



The save function is activated when the CBserver option {\tt -save} {\it savedir\/} is specified.
The source files are saved in the directory specified by the {\it savedir} parameter.
This parameter must be the path to an existing directory. If activated, the
save function is invoked upon the following events:

\begin{enumerate}
\item The CBserver is shutdown. In this case, the complete module tree starting from
root module {\tt System} is saved.
\item A client disconnects from the CBserver. Here, the module tree starting at the home
module of the user associated to the client tool is saved.
\item A user changes the module context. In this case, the old module is saved.
\end{enumerate}

In all cases, the save function is executed with the
rights of the user who started the CBserver. The save function requires at least read
permission for the module to be saved. The above rules are also applicable to 
the server-side materialization of query results, see section \ref{sec:module_views}.

The load function gets activated when the option {\tt -load} {\it loaddir\/} is specified.
The directory {\it loaddir\/} should contain files with file/directory names being formed
as explained for the save function. It loads the files in {\em alphabetic order} to control
the sequence in which the files are loaded. If you manually add Telos source files
to a directory that is about to be loaded, then make sure that its file name is alphabetically
sorted after the module file name, e.g. {\tt AB\_01extension.sml} is loaded after
{\tt AB.sml}.
The import is executed once at CBserver startup with the rights of the user who
started the CBserver. If a file contains an error, the loading of this module source
fails. Error messages shall be displayed in the trace log of the CBserver. Note
that the CBserver can be started with a non-empty database. The import of source files
will be added to the already existing content of the database.

Examples:

\begin{verbatim}
   cbserver -d DB1 -save /home/meee/DB1SRC
\end{verbatim}

This command starts up a CBserver that will eventually save the module sources in the specified
 directory. Note that the saving takes place either at CBserver shutdown or when a
client tool disconnects (partial save).

\begin{verbatim}
   cbserver -u nonpersistent -save /home/meee/SRC
\end{verbatim}

This variant will start a CBserver with a non-persistent database but the contents will nevertheless
stored as Telos sources file in {\tt /home/meee/SRC}.


\begin{verbatim}
   cbserver -u nonpersistent -load SRC1 -save SRC2
\end{verbatim}

This command starts a the CBserver (i.e. only system objects are defined) and then
loads module sources from the directory {\tt SRC1}. Then, client tools may
modify the contents of the database. Finally, the module contents are saved in 
directory {\tt SRC2}.

\begin{verbatim}
   cbserver -d DB1 -load /home/meee/DB1SRC -save /home/meee/DB1SRC
\end{verbatim}

The load and save directories may also be the same. Note that the save function
will eventually overwrite the files that have been loaded at CBserver startup.

\begin{verbatim}
   cbserver -u nonpersistent -load DB1SRC 
\end{verbatim}

This command will start a non-persistent CBserver and loads the module sources of
directory DB1SRC. 


Module sources do not contain the historic states of objects that are
maintained with rollback times in the CBserver database. Hence, a persistent database
is containing more information than the saved module tree and is also faster to
start up compared to loading module sources. Still, the load/save function offers
a simple way to keep a textual representation synchronized with the 
evolving database state, or to back-up/re-load a database state. 
The CBserver parameter {\tt -db}  combines the function of {\tt -d}, {\tt -load}, 
{\tt -save}, and {\tt -views}. Hence, all files will be accessed/updated in
the database directory.






\section{Server-side materialization of query results} \label{sec:module_views}
% written by M.Jeusfeld, 22-Feb-2011

Similar to the saving of module sources, the CBserver parameter {\tt -views}
enables the materialization of certain query results in the file system of
the CBserver.
To do so, one has to specify the queries to be materialized. Only queries with a
single parameter or with no parameter can be materialized. The queries need to be listed with the
module that contains the objects that match the query. Example:

\begin{verbatim}
   MyModule with
     saveView
       v1: Q1;
       v2: Q2
   end
\end{verbatim}

You can also use deductive rules deriving the values for the {\tt saveView} attribute.

The queries {\tt Q1} and {\tt Q2} need to be visible in the module {\tt MyModule}.
The queries need to have an answer format (section \ref{sec:answerformat})
defined for them (attribute {\tt forQuery}).
Assume that the query {\tt Q1}  has the single parameter {\tt param:C1}. ConceptBase
will then call the query {\tt Q1[x/param]} for each instance {\tt x}  of class {\tt C1}.

The result is stored in a file with name {\tt x} in the directory specified with
the {\tt -views} parameter. If the view is extracted from a module different to
{\tt oHome}, then the filename includes the module name as a prefix. 
The file type of the file is taken from the optional
{\tt fileType} attribute of the answer format of {\tt Q1}. The default file type
is ".txt".
The result of queries with no parameter is stored in files carrying the name of the
query.
If the module separator is set to '/', then the files are stored in sub-directories
that reflect the module path, from which the view was extracted, see
also section \ref{sec:module_sources}.

The materialization of query results is executed at the same events when the 
saving of module sources takes place (section \ref{sec:module_sources}).
To enable the materialization, you need to 
specify the target directory with the {\tt -views} option:

\begin{verbatim}
   cbserver -d MYDB -views /home/meee/MyViews 
\end{verbatim}

You can also use the CBShell utility (section \ref{cha:cbshell}) to extract
the query results and materialize them on the {\em client side}. This method
is more flexible but you need to program the CBShell scripts for
to extract all required views. For example, the CBShell script

\begin{verbatim}
   connect alpha 4001 
   ask Q1[x/param] OBJNAMES default Now
   showAnswer
   exit
\end{verbatim}

connects to the CBserver running on a host named {\tt alpha} with port number
4001 and will extract just the answer to {\tt Q1[x/param]}. If there are more answers to
be extracted, one has to employ separate scripts for each of them and execute them
one after the other to save the results in separate files.
The {\em server-side\/} method using the {\tt -views} option will determine
all possible fillers {\tt x} for the parameter {\tt param} and automatically save
the results of {\tt Q1[x/param]} in a separate file with filename {\tt x}.

You can use the {\tt -db} option for activating the saving/loading of module sources
and the materialization of query results within the database directory:

\begin{verbatim}
   cbserver -db MYDB 
\end{verbatim}

This creates a single directory with both the binary database files, the module sources,
and the materialized query results.
Further examples are available from the CB-Forum at \url{http://merkur.informatik.rwth-aachen.de/pub/bscw.cgi/3097259}.


\subsection{Post-export command}

ConceptBase can only export textual query results. Some output formats such as
program source code can be further processed by calling appropriate tools.
This can be initiated manually, or one can configure the directory specified
in the {\tt -views} option with a command file  {\tt postExport.sh}.
This command file shall then be executed
by the CBserver whenever a saving operation on the views directory has taken place.

You have to regard that the command is executed in the context of the CBserver.
Hence, it is executed in the directory in which the CBserver was started.
You should therefore change to the views directory inside the post-export
command file like shown in this example of {\tt postExport.sh}:

\begin{verbatim}
   #!/bin/sh
   cd `dirname $0`
   cp *.xml /home/meee/public
\end{verbatim}

The second line changes to the views directory. The third line does the 
specific processing of the materialized query results. Here, we just copy the
file. More interesting are calls to transformation routines. 

You should remove the write permission for the post-export command to prevent that
it can be overwritten%
\footnote{Overwriting the post-export command file may be a desired feature. You can then
generate it from within the CBserver like any other file in the views directory. 
This is however a major security hole and we strongly discourage to use this feature.}
by the materialization function. Under Unix/Linux, this can be achieved via the command

\begin{verbatim}
   chmod u-w postExport.sh
\end{verbatim}

The above command only has to be executed once within the views directory.
If you use the module separator '/' (i.e. the deep structure explained in section \ref{sec:module_sources}),
then the view files are stored in the directory of the module they belong to.
This also may effect the way how you program the postExport script file. In Unix/Linux,
you can use the 'find' command to fetch all files that are subject for post-processing:

\begin{verbatim}
   #!/bin/sh
   cd `dirname $0`
   xmlfiles=`find . -name "*.xml"`
   for file in $xmlfiles; do
    cp $file /home/meee/public
   done
\end{verbatim}

This script also works fine in the case of a flat view directory structure.


